\scanpage{21}
% archive sheet number: 20
\unsure{non} l'égalité -- celle-ci sert à \unsure{sommer} (72).

La situation devient encore plus jolie si on suppose
$\unsure{\zeta|z}=1$, puisque alors la donnée supplémentaire
revient à la donnée
(\unsure{69|09|0g|eq})\qquad un sous-groupe \inserted{(commutatif)} $T_H$ de $H$, normalisant $U_0$,
tel que $T_H\cap U_0=\{1\}$, et que $\sigma(t)=t^{-1}$ pour
$t\in T_H$.
% A long diagonal runs across the square and lower working, ending at the last braid word before the two final tests.
\deleted{\diagram{\begin{tikzcd}[ampersand replacement=\&]
\widetilde Z\arrow[r,"\unsure{k|K|\varepsilon}"]\arrow[d,"u_0"'," \unsure{2|\wr}"] \& Z\arrow[d,"u"," \unsure{2|\wr}"']\\
T_0\arrow[r,"c_0"] \& T_H
\end{tikzcd}}
\diagram{$\begin{array}{cc}\lambda^{-1}&\lambda^{\unsure{\omega|m}}\\1&\lambda^2\end{array}$}
\[\widetilde Z\]

$\varepsilon_0,\varepsilon_1,\sigma,\rho_0,\rho_1,\rho'_0,\rho'_1,\omega$
\margin{$\rho_0^3=\rho_1^3=\sigma^2=\omega$\\ $\omega^2=1$, $\omega$ commute à tout}
\[\left\{\begin{aligned}
\rho'_0&=\sigma \varepsilon_0 &\rho_0&=\inserted{\sigma^{-1}\varepsilon_0=}(\varepsilon_1\varepsilon_0)^{-1} &\rho'_0&=\inserted{\deleted{\unsure{\omega\rho\varepsilon_0}}}\rho_0^{-2} &{\rho'_0}^3&=1 &\rho_0^6&=1\\
\rho'_1&=\sigma \varepsilon_1 &\rho_1&=\inserted{\sigma^{-1}\varepsilon_1=}(\varepsilon_0\varepsilon_1)^{-1} &\rho'_1&=\rho_1^{-2} &{\rho'_1}^3&=1 &\rho_1^6&=1
\end{aligned}\right.\]
\margin{$\rho'_0=\omega\rho_0,\ \rho_0=\omega\rho'_0$\\ $\rho'_1=\omega\rho_1,\ \rho_1=\omega$}
\[\left\{\begin{aligned}
\varepsilon_1&=\sigma(\varepsilon_0)=\rho_{\unsure{1|0}}(\varepsilon_0)=\rho'_0(\varepsilon_0)=\varepsilon_1(\varepsilon_0)\\
\varepsilon_0&=\sigma(\varepsilon_1)=\rho_0(\varepsilon_1)=\rho'_1(\varepsilon_1)=\varepsilon_0(\varepsilon_1)\\
\sigma&=\varepsilon_1\varepsilon_0\varepsilon_1=\varepsilon_0\varepsilon_1\varepsilon_0
\end{aligned}\right.\tag*{$(\ast)$}\]

Partons de $\varepsilon_0,\varepsilon_1$ [satisfaisant $(1)\ \boxed{\varepsilon_1\varepsilon_0\varepsilon_1=\varepsilon_0\varepsilon_1\varepsilon_0}$]
posons $\sigma=\varepsilon_0\varepsilon_1\varepsilon_0[=\varepsilon_1\varepsilon_0\varepsilon_1],\rho'_0=\sigma \varepsilon_0,\rho'_1=\sigma \varepsilon_1,\rho_0^{\unsure{?|-1}}=\varepsilon_1\varepsilon_0,\rho_1^{\unsure{1|-1}}=\varepsilon_0\varepsilon_1$
\deleted{je dis qu'on a bien} \inserted{examinons les formules $(\ast)$} \deleted{\unsure{alors}} qui restent à vérifier. On a
\[\underbrace{\sigma(\varepsilon_0)=\rho'_0(\varepsilon_0)}=\varepsilon_0\varepsilon_1(\varepsilon_0)\deleted{\unsure{\varepsilon_0\varepsilon_0^{-1}\varepsilon_1^{-1}\varepsilon_1^{-1}}}=\varepsilon_0\varepsilon_1\varepsilon_0\varepsilon_1^{-1}\varepsilon_0^{-1}\stackrel{?}{=}\varepsilon_1\]
c'est égal à $\varepsilon_0$ ssi \deleted{$\varepsilon_1\varepsilon_0=\varepsilon_0\varepsilon_1$} on a (1)
\[\sigma(\varepsilon_1)=\rho'_1(\varepsilon_1)=\varepsilon_0\underline{\varepsilon_1\varepsilon_0\varepsilon_1}\varepsilon_0^{-1}\varepsilon_1^{-1}\varepsilon_0^{-1}\stackrel{?}{=}\varepsilon_0\]
égal à $\varepsilon_0$ ssi $\varepsilon_1\varepsilon_0\varepsilon_1=\varepsilon_0\varepsilon_1\varepsilon_0$ i.e (1)}
\[\begin{aligned}
\rho_0^2&=\varepsilon_1\varepsilon_0\varepsilon_1\varepsilon_0\stackrel{?}{=}\rho'_0=\varepsilon_0\varepsilon_1\varepsilon_0\mathbin{\cdot}\varepsilon_0\quad\text{ssi on a (1)}\\
\rho_1^2&=\varepsilon_0\varepsilon_1\varepsilon_0\varepsilon_1\stackrel{?}{=}\rho'_1=\varepsilon_0\varepsilon_1\varepsilon_0\varepsilon_1\quad\text{tjrs vrai}
\end{aligned}\]
